\section{Other branching numbers and the limit of the small anchor}
\label{sec:branching}

\subsection{A family on regular rooted trees}

The full-pattern pinning theorem separates the size of the anchor from
the tree arithmetic. This gives a uniform extension to every fixed
branching number.

\begin{theorem}[Regular branching]
\label{thm:branching}
Fix an integer \(b\ge2\), put \(s=b+1\), and let
\(H^{(b)}=K_{b+1}(P)\) with the body \(P\) in \eqref{eq:H}.
Thus \(H^{(b)}\) has size \((b+3)\times(b+3)\). For every \(h\ge1\),
there is an explicit binary matrix \(A_{b,h}\) with
\begin{equation}
 m=b^h,\quad D_b=2b^2+4b+4,\quad d_h=D_bh+2,\quad n=d_hm,
 \label{eq:bary-size}
\end{equation}
such that
\begin{equation}
 N_{H^{(b)}}(A_{b,h})=b\,m^{2b+4},\qquad
 \dist_{H^{(b)}}(A_{b,h})=d_h^{-2},\qquad
 \tau_{H^{(b)}}(A_{b,h})=\nu_{H^{(b)}}(A_{b,h})=m.
 \label{eq:bary-exact}
\end{equation}
There is an explicit optimal repair of \(m^2\) cells. For independent
uniform balanced samples with \(q/(d_hm^{2/3})\to c\in(0,\infty)\),
the leaf-witness count converges to \(\Poi(bc^3)\), and detection
probability tends to \(1-e^{-bc^3}\).
\end{theorem}
\begin{proof}
Replace the binary tree by the complete \(b\)-ary tree. A block at
depth \(i\) has size \(m/b^i\); nodes are indexed from zero to
\(b^i-1\), and \(\pi(p)=\lfloor p/b\rfloor\). Use the same
recursive orders, placing all \(b\) child subtrees in increasing
order between the two signed blocks. In \eqref{eq:roles}, retain
the two \(V\) modes and use \(W_{i,r}^+,W_{i,r}^-\) for each
residue \(r\in\{0,\ldots,b-1\}\). The same mode order is used
on both axes. Use \(s=b+1\) anchor groups for each mode, with the
same crown and unit-vector signatures. All variable entries are
given by \eqref{eq:entries} with the new parent map.

There are \((2+2b)h\) modes, \(2h+1\) variable classes, and one
dummy class per axis. Their total gives
\[
 d_h=((2+2b)(b+1)+2)h+2=(2b^2+4b+4)h+2.
\]
Every full variable class still has size \(m\). Threshold traces
remain nested, and the anchor signature of a member varies only
with its residue modulo \(b\). Thus there are at most \(b=s-1\)
signatures. The support-pair graph is still two signed paths joined
by the terminal triangle. Its leaf traces remain a chain, now with
\(\lfloor p/b\rfloor\) in place of \(\lfloor p/2\rfloor\).
The crown-exclusion proof, including after the dummy-by-leaf repair,
therefore applies verbatim. Every hypothesis of \cref{thm:pinning}
is verified.

The four order comparisons remain valid with the new parent map.
The body classification uses only those comparisons and the fact
that two child indices in the same residue class have the same
relative order as their parents. It therefore excludes all modes
except \(W_{h,p\bmod b}^-\), where the leaf row and column must
have the same index and the other column belongs to their parent.
That parent block now has \(b\) columns. There are \(m\) leaves,
\(m\) dummy rows, and \(m^{2s}\) anchor choices, giving
\(b m^{2s+2}=b m^{2b+4}\) copies.

The leaf diagonal is an \(m\)-cell transversal. To obtain \(m\)
copies with disjoint row and column sets, assign the \(b\) columns
of each parent bijectively to its children and use distinct dummy
and anchor representatives, as in \cref{thm:packing}. For the edit
lower bound, a uniform leaf chooses its depth-\(i\) ancestor with
probability \(b^{-i}\), so a uniform representative in that node's
block is again uniform on its size-\(m\) class. All protected-cell
probabilities are \(m^{-2}\). The same four propagation bodies
force opposite values at the shared leaf. The explicit repair
\(R_*\times C_h\) excludes every pinned final mode, proving the
matching upper bound.

For the sampling claim, let \(p=q/n\) and let \(X\) count leaves
whose row, matching column, and at least one of their \(b\) parent
columns are sampled. At Bernoulli retention probability \(p\), its
exact generating function is
\begin{equation}
 \E z^X=\left[1-t_b+t_b(1-p^2+p^2z)^b\right]^{m/b},
 \qquad t_b=1-(1-p)^b.
 \label{eq:bary-pgf}
\end{equation}
For fixed-size samples, expand the \(r\)-th factorial moment by
the number \(j\) of distinct parents of its \(r\) leaves.
The distinct-parent tuples number \(m^r(1+O_{b,r}(m^{-1}))\).
Their joint probability is
\(b^r p^{3r}(1+O_{b,r}(p+q^{-1}))\), by inclusion-exclusion on
the disjoint parent blocks. Their contribution tends to
\((bc^3)^r\). Tuples with \(j<r\) parents number \(O_{b,r}(m^j)\)
and have probability \(O_{b,r}(p^{2r+j})\); their contribution is
\(O_{b,r}(m^{-2(r-j)/3})=o(1)\). The earlier factorial-moment
argument gives \(X\Rightarrow\Poi(bc^3)\). All \(b\) final
minus-mode anchors and the dummy row are covered except with
probability at most \((2bs+1)e^{-q/d_h}=o(1)\), proving the
detection limit.
\end{proof}

Put \(a_b=\log b/D_b\). The exact copy density for this family is
\[
 \frac{N_{H^{(b)}}(A_{b,h})}{n^{2b+6}}
 =b\,d_h^{-2b-6}m^{-2}
 =b e^{4a_b}\eps_h^{b+3}
       \exp(-2a_b\eps_h^{-1/2}).
\]
For fixed \(b\), the exponential coefficient is \(2a_b\).

\begin{proposition}[Binary branching optimizes this family]
\label{prop:binary-optimal}
Among integers \(b\ge2\), \(a_b=\log b/(2b^2+4b+4)\) has its
unique maximum at \(b=2\).
\end{proposition}
\begin{proof}
For real \(x\ge2\), the derivative of
\(g(x)=\log x/(2x^2+4x+4)\) has the sign of
\[
 f(x)=2x+4+4/x-(4x+4)\log x.
\]
Its derivative is \(-2-4/x-4/x^2-4\log x<0\), and
\(f(3)=34/3-16\log3<0\). Thus \(g\) decreases on \([3,\infty)\).
The remaining comparison \(g(2)>g(3)\) is equivalent to
\(17\log2>10\log3\), which follows from
\(2^{17}=131072>59049=3^{10}\).
\end{proof}

This optimization concerns the displayed regular-tree family. It does
not constrain constructions using different block weights, variable
branching, different signatures, or another propagation mechanism.

\subsection{Why simply removing another anchor group fails}

The proof of \cref{thm:pinning} uses the body positions essentially;
the three-by-three crown alone is not rigid. Decreasing to two anchor
groups would give the pattern
\begin{equation}
 H_4=\begin{pmatrix}
 0&1&1&0\\1&0&0&1\\1&0&1&0\\0&1&1&1
 \end{pmatrix}.
 \label{eq:H4}
\end{equation}
The same host and repair no longer work.

\begin{proposition}[A four-by-four copy surviving the proposed repair]
\label{prop:H4-counterexample}
Form the binary host at depth \(h=2\) using \(s=2\) anchor groups.
Both this host and its dummy-by-leaf repair contain an unpinned copy
of \(H_4\).
\end{proposition}
\begin{proof}
Select, in order, the following four rows:
\[
 \widehat R_{V_2^+}(2),\quad
 \widehat R_{W_{2,0}^-}(2),\quad R_2(2),\quad R_*.
\]
One arbitrary representative is chosen from each displayed block.
Select columns from
\[
 \widehat C_{W_{2,0}^-}(1),\quad
 C_1^+(0),\quad C_1^+(1),\quad C_1^-(1).
\]
The mode order gives the row order, and the recursive column order
gives the stated column order. The first two rows have traces
\(0110\) and \(1001\) by their anchor incidences. For leaf index
\(2\), the parent index is \(1\), so its trace is \(1010\).
The dummy row has trace \(0111\). These are exactly the rows of
\(H_4\). There are no selected leaf columns, so the repair changes
none of these entries. Only one selected column is an anchor,
exhibiting the failure of full-pattern pinning.
\end{proof}

This is a counterexample to this particular two-group construction
and repair argument. It neither proves polynomial removal for every
four-by-four pattern nor excludes a different four-by-four obstruction.

\section{Research questions suggested by the exact construction}
\label{sec:questions}

The preceding theorems settle concrete refinements of the reviewed
source construction. The following questions identify further work;
none is used as an assumption in the proofs.

\begin{question}[Small forbidden matrices]
What is the least \(k\) for which a binary \(k\times k\) pattern
fails polynomial ordered removal? The present work gives \(k\le5\).
Determine which patterns of orders two, three, and four admit
polynomial removal. In particular, can one control the additional
\(H_4\)-copies in \cref{prop:H4-counterexample} by a different repair
and counting argument, or find a different small anchor-body pattern?
\end{question}

\begin{question}[The strength of the removal obstruction]
For the explicit \(H\) in \eqref{eq:H}, determine the asymptotic
order of \(\mathcal R_H(\eps)\). Can the construction be improved
to give \(\mathcal R_H(\eps)\le\exp(-c\eps^{-\alpha})\) with
\(\alpha>1/2\), or can a matching lower bound of stretched-exponential
type be proved? Reducing the number of protected classes, allowing
unequal block sizes, or sharing anchors across levels are concrete
directions. \Cref{prop:binary-optimal} rules out an improvement of
the exponential coefficient from changing only the regular branching
number in the displayed family.
\end{question}

\begin{question}[Universal repair from a transversal]
Define
\[
 \Gamma_H(t)=\sup\left\{\dist_H(A):
         \tau_H(A)/n^2\le t,\quad A\in\{0,1\}^{n\times n},\ n\ge1\right\}.
\]
Does \(\Gamma_H(t)=O(\log^{-2}(1/t))\) hold for the five-by-five
pattern, or can the gap between original-copy covering and binary
repair be larger? The construction proves
\(\liminf_{t\downarrow0}\Gamma_H(t)\log^2(1/t)\ge a^2\).
A matching upper bound would turn the inverse logarithmic obstruction
into a sharp structural theorem for all hosts.
\end{question}

\begin{question}[Structure and stability of optimal repairs]
Classify all \(m^2\)-cell repairs of \(A_h\). More generally, what
must an \(H\)-free matrix differing in at most \((1+\eta)m^2\)
cells look like as \(\eta\downarrow0\)? Is there a finite collection
of repair mechanisms near which every near-minimizer lies, or do
many distributed repairs attain the same cost? Exact distance alone
does not resolve the geometry of the minimizing edits.
\end{question}

\begin{question}[A quantitative Poisson approximation]
Obtain explicit total-variation error bounds for the fixed-size
witness law, uniformly when \(q/(d_hm^{2/3})\) ranges in a compact
subset of \((0,\infty)\). Separate the correction to the mean,
sibling dependence, and sampling without replacement. Determine the
first finite-depth correction to the detection quantile, including
the error from anchor coverage. The exact factorial moments and
Bernoulli generating function provide independent starting points.
\end{question}

\begin{question}[Uniqueness and asymptotics of the optimal rectangle]
Is \(\theta\mapsto\Phi_s(z,\theta)\) unimodal for every \(s\ge3\)
and \(z>0\)? Determine the optimizer and \(z_{\rho,s}\) as
\(\rho\downarrow0\), \(\rho\uparrow1\), or \(s\to\infty\).
The interior-optimizer theorem does not imply uniqueness. It would
also be useful to determine whether randomized sample budgets or
several rectangles improve expected entry cost at a fixed detection
probability.
\end{question}

\begin{question}[Adaptive information]
What lower bounds survive when a tester may choose arbitrary entries
adaptively? A fixed explicit host has copies at known locations, so
uniform-sampling lower bounds alone cannot answer this. Construct
distributions of ordered hosts that hide the relevant modes and leaf
witnesses while preserving distance, and analyze one-sided and
two-sided adaptive testers separately. The extension argument in
\cref{prop:completion} suggests a precise information-completion
criterion for partial observations beyond rectangles.
\end{question}

\begin{question}[Other witness exponents]
Can related ordered constructions force a critical witness to use
\(r\) rare sampled positions, for fixed \(r>3\), while retaining
an exact repair theorem and a fixed forbidden pattern? With \(m\)
potential witnesses, the heuristic scale would become
\(q\asymp d_hm^{1-1/r}\). A proof must control overlapping witnesses
and anchor coverage; the heuristic by itself supplies neither.
\end{question}

\begin{question}[General full-pattern pinning]
Replace the crown and the two unit-vector signatures in
\cref{thm:pinning} by other fixed binary patterns. Which combinatorial
conditions on nested classes, signature counts, and the body force
the prefix to lie in a single anchor mode? A classification could
separate the part of an anchor that is rigid by itself from the
additional rigidity supplied by the body, and provide smaller
obstructions in other ordered induced-removal problems.
\end{question}

\section{Source record, verification, and integration}
\label{sec:verification}

\subsection{What was reviewed and what is proved here}

The source review is fixed to the following repository snapshots,
retrieved on 7 October 2026:
\begin{itemize}
\item \texttt{openai/math}, commit\\
\texttt{adc7f1241b42e322a6451854ab7e4b4c146bf78a}.
The selected item is Family 190, \emph{Polynomial removal fails for
ordered binary matrices}, manuscript date 25 September 2026.
The construction, original anchor argument, distance proof, counting
argument, conclusion, and bibliography were read. Its formalization
catalogue identifies a Lean comparator for the \(66\times66\)
nonpolynomial-removal theorem~\cite{OAI190Lean}; that upstream project
was not rebuilt during this work.
\item \texttt{VladimirReshetnikov/ProveIt}, commit\\
\texttt{fb9602e556560a15704b4feb176ba6de08378665}.
The relevant sources are the transseries theorem
\texttt{p0:thm:lambert-core} and the file
\texttt{LinLogCoreInversion.lean}. The exact positive-parameter
inversion and uniqueness statements were checked at the source level.
Those existing Lean modules were not rebuilt here.
\end{itemize}

OpenAI's repository explicitly distinguishes results at different
stages of verification~\cite{OAIRepo}. ProveIt's repository similarly
distinguishes formal theorems, executable certificates, and exploratory
research material~\cite{ProveIt}. These distinctions motivate the
choice of a finite combinatorial construction whose essential steps
can be written out and tested independently. This article is a focused
research extension, not a certification of all results in either
repository.

The inherited ingredients are the signed-tree host, its role table and
orders, its nonanchor threshold entries, and its protected-path
propagation lower bound. The new refinements established here are the
full-pattern pinning theorem and five-by-five obstruction; the repair
that attains the lower bound and the proof that it creates no copies;
the exact classification, count, packing, and transversal; the finite
sampling identities and sharp limits; the optimal uniform-rectangle
cost; the exact transversal-to-repair relation with its explicit
convergent logarithmic expansion; and the stated regular-tree
generalization. The standard Lambert inversion is credited separately.
Novelty is asserted relative to the reviewed source results; no claim
of a global minimum pattern size or an exhaustive priority search is
made.

\subsection{Reproducible finite and numerical checks}

The companion program \texttt{code/verify.py} uses only Python's
standard library, exact integers, and rational arithmetic. It implements
the construction independently of the written pinning proof. Equal
rows or columns inside a block are represented by one weighted
position. Since every row and every column of \(H\) is distinct,
an \(H\)-copy cannot choose two positions from one such twin block;
weighted compressed enumeration therefore counts the expanded host
exactly.

The default checks include the following:
\begin{itemize}
\item Exhaustive searches for complete \(H\)-copies at depths one and
two, in both the original and repaired hosts. The compressed orders
are \(23\) and \(47\); the search examines all \(33{,}649\) and
\(1{,}533{,}939\) row five-tuples, respectively, and all compatible
column chains. Every original copy has its full prefix pinned.
The exact expanded counts are \(512=2\cdot2^8\) and
\(131{,}072=2\cdot4^8\); both repaired counts are zero.
\item Structural nesting, signature, and nonanchor crown checks through
depth six; explicit disjoint expanded copies; and weighted body
enumeration in every designated mode. Completeness of the latter
enumeration for arbitrary depth is justified by the written pinning
theorem, not inferred from testing.
\item An explicit \(H_4\)-copy surviving repair, as in
\cref{prop:H4-counterexample}, and small-host checks of the failure
of crown-only pinning.
\item Exact Bernoulli formulas checked by enumerating all relevant
leaf-row, leaf-column, and parent-column outcomes in small instances.
The independent anchor and dummy events are integrated analytically.
\item The fixed-size factorial-moment identity checked against every
row and column sample in witness submodels \((m,n)=(2,6),(4,10)\),
for all \(q\in\{0,\ldots,n\}\) and every order \(1\le r\le m\):
\(58\) exact rational comparisons in total.
\end{itemize}

The output \texttt{code/verification\_results.json} records the
parameters and outcomes. The checks use explicit exceptions, so they
also run under Python's optimization flag. They check the repair's
size and small-depth freeness; they do not search all possible edit
sets or computationally prove the general \(m^2\) lower bound.

The optional script \texttt{code/verify\_inversion.py} checks the
first logarithmic coefficient polynomials symbolically when SymPy is
installed, and evaluates
the exact inverse and remainder estimates at high precision.
\texttt{code/generate\_figures.py} reproduces the vector figures and
the illustrative rectangular constants using NumPy, SciPy, and
Matplotlib. Its probabilities are evaluations of exact formulas;
the numerical optimizers are approximations, not certified interval
enclosures. The analytic theorems do not rely on those optimizations.

\subsection{A concrete route to formalization}

The finite structural core offers a natural first integration target
for ProveIt. A staged dependency order is preferable to treating
the whole article as one formalization obligation:
\begin{enumerate}
\item Define ordered binary copies by increasing embeddings and prove
the twin-block compression lemma and the completion lemma.
\item Formalize \cref{thm:pinning} as an abstract finite theorem.
Its proof uses order, distinct traces, set inclusion, and a bounded
number of signatures; it does not require probability or limits.
\item Define the signed tree positions and entries, then prove the
four order comparisons, the signature bound, and nonanchor crown
exclusion for the original and repaired hosts.
\item Prove copy classification, exact count, the explicit packing,
and repair freeness. These give finite identities for every \(h\).
\item Formalize the arbitrary-edit lower bound by counting protected
selections, or by finite probability. The shared-path dependence
must remain explicit; the proof only uses independence for protected
cells, not for arbitrary variable pairs.
\item Add finite sampling identities. Poisson convergence and
rectangular asymptotics can follow as separate analytic modules.
For the inverse, reuse ProveIt's existing principal Lambert core
and formalize the specialized combinatorial substitution and the
analytic contraction remainder.
\end{enumerate}

The archive supplies the article, complete LaTeX build inputs, exact
verification code, numerical scripts, and their receipts. It introduces
no new axiom declarations or unproved Lean stubs. The new theorems are
presented as mathematical proofs with supplementary checks, ready for
independent review and subsequent formalization.
